\title{xLSTM: Extended Long Short-Term Memory}

\begin{abstract}
The foundational concepts of the Long Short-Term Memory (LSTM) architecture—namely, the constant error carousel and gating mechanisms—emerged in the 1990s and have since become pivotal in the evolution of deep learning, forming the backbone of the initial Large Language Models (LLMs). Despite their longevity and impact, LSTMs were eventually surpassed in scalability and performance by the Transformer architecture, which is characterized by highly parallelizable self-attention mechanisms. This shift prompts a fundamental inquiry: by scaling LSTMs to the order of billions of parameters, incorporating state-of-the-art techniques from contemporary LLMs, and systematically addressing traditional LSTM shortcomings, how effectively can LSTMs serve language modeling tasks today? To address this, we first propose exponential gating, augmented by suitable normalization and stabilization procedures. We also introduce restructured LSTM memory mechanisms: (i) the sLSTM variant features a scalar memory and scalar update, paired with a novel memory mixing strategy; (ii) the mLSTM version is equipped with a matrix memory and leverages a parallelizable, covariance-based update scheme. These enhanced LSTM designs are integrated into residual block frameworks, producing xLSTM units, which are then hierarchically organized into comprehensive xLSTM models. The combined improvements in gating and memory architecture significantly strengthen xLSTM performance, enabling them to match or surpass the latest Transformer and State Space Model benchmarks in both scaling properties and task effectiveness.\\
Code available at: \url{https://github.com/NX-AI/xlstm}
\end{abstract}

\section{Introduction}
\label{sec:intro}

The concepts behind Long Short-Term Memory (LSTM)~\citep{Hochreiter:91,Hochreiter:97nips,Hochreiter:97}, specifically the gating mechanisms and the constant error carousel, were proposed as solutions to address the vanishing gradient issue in recurrent neural networks~\citep{Hochreiter:91,Hochreiter:00book}:

\begin{align}
\label{eq:lstm_idea}
  c_t \ = \ f_t \ c_{t-1} \ + \ 
          i_t \  z_t \ , \quad  
          h_t \ = \ o_t \ \psi({c_t}) \ . 
\end{align}

The constant error carousel refers to the process by which the cell state $c_{t-1}$ (depicted in green) is incrementally adjusted through the addition of cell inputs $z_t$, with this modification being regulated via sigmoid gates (shown in blue). Specifically, the input gate $i_t$ and the forget gate $f_t$ govern these cell state updates, whereas the output gate $o_t$ determines the extent to which the memory cell’s content is reflected in the output, that is, the hidden state $h_t$. Before producing the hidden state, the cell state undergoes normalization or squashing through the function $\psi$, after which output gating yields the hidden state.

LSTMs have achieved notable success across a diverse array of fields \citep{Hochreiter:01,Hochreiter:07,Schmidhuber:15}, dominating text generation tasks prior to the widespread adoption of Transformers in 2017~\citep{Vaswani:17}. Their strong performance has been validated in various sequential problem domains including, but not limited to, text generation~\citep{Graves:13,Karpathy:15blog}, handwriting synthesis~\citep{Graves:13}, machine translation using sequence-to-sequence architectures~\citep{Sutskever:14nips}, program evaluation~\citep{Zaremba:14arxiv}, automated image captioning~\citep{Karpathy:15,Hossain:19}, source code generation~\citep{Karpathy:15blog}, hydrological modeling of rainfall-runoff processes~\citep{Kratzert:18,Kratzert:19}, and development of flood early warning systems~\citep{Nearing:24}. Within reinforcement learning, LSTMs have consistently ranked as top-performing sequence models in prominent systems such as AlphaStar for StarCraft II~\citep{Vinyals:17short}, OpenAI Five for Dota 2~\citep{OpenAI:19}, and models steering magnetic controllers in nuclear fusion experiments~\citep{Degrave:22short}. A key strength of LSTMs lies in their ability to learn high-level representations by efficiently capturing and retaining semantic content within their memory mechanisms~\citep{Karpathy:15blog}, as reflected in the discovery of specialized neurons associated with numerical reasoning and syntax~\citep{Lakretz:19}, linguistic attributes~\citep{Bau:19}, and sentiment detection~\citep{Radford:17arxiv}. Even today, LSTMs are integral to state-of-the-art applications~\citep{Degrave:22short,Nearing:24}, underscoring their enduring impact and resilience.

Although LSTMs have achieved significant breakthroughs, they are hampered by three primary limitations: (i) Lack of ability to update stored information. We illustrate this constraint using the {\it Nearest Neighbor Search} problem (additional discussion in Appendix~\ref{sec:appNearestNeighborSearch}): When provided with a reference vector, the task involves traversing a sequence to identify and retrieve the value associated with the most similar vector at the end of the sequence. As depicted in the left panel of Figure~\ref{fig:lstmProblems}, LSTMs exhibit elevated mean squared error on this task; specifically, they cannot overwrite a previously stored value should a more similar candidate emerge later in the sequence. In contrast, our proposed xLSTM overcomes this issue through the use of exponential gating. (ii) Constrained memory capacity, whereby all information must be encoded into scalar-valued cell states. This shortcoming is highlighted in the context of {\it Rare Token Prediction}. In the right panel of Figure~\ref{fig:lstmProblems}, perplexity results are reported on Wikitext-103~\citep{Merity:17}, stratified by token frequency. LSTM models show substantially higher perplexity on rare tokens due to the limitations of their storage mechanism, whereas xLSTM employs a matrix-based memory to address this shortfall. (iii) Inefficiency in parallel processing as a result of memory entanglement; the hidden-to-hidden connections across time steps necessitate inherently sequential operations.

The constraints inherent to LSTM have contributed to the rise of Transformers~\citep{Vaswani:17} in the field of language modeling. By addressing these shortcomings and scaling LSTM architectures to match the magnitude of contemporary Large Language Models, an important question emerges: what levels of language modeling performance become attainable?

\section{Extended Long Short-Term Memory}
\label{sec:xLSTM}

To address the shortcomings of LSTM, Extended Long Short-Term Memory (xLSTM) makes two primary advances over the foundational concept shown in Equation~\eqref{eq:lstm_idea}. These advances—namely, exponential gating and innovative memory architectures—broaden the LSTM framework by adding two distinctive variants: (i) sLSTM (described in Section~\ref{sec:sLSTM}), which employs a single scalar as its memory, performs scalar updates, and utilizes memory mixing; and (ii) mLSTM (outlined in Section~\ref{sec:mLSTM}), which incorporates a matrix-based memory and adopts a covariance update rule grounded in outer products, allowing for full parallelization. Exponential gating serves as a unifying enhancement in both sLSTM and mLSTM. To support parallel execution, the mLSTM omits memory mixing, effectively removing recurrent connections between hidden states. Both mLSTM and sLSTM architectures permit the expansion to multiple memory cells, with sLSTM leveraging memory mixing across different cells. Additionally, sLSTM can include multiple heads—each head operates without memory mixing between them, while still enabling mixing among the cells contained within each head. The incorporation of heads alongside exponential gating in sLSTM creates an alternative mechanism for memory mixing. In contrast, for mLSTM, there is no distinction between the concepts of multiple heads and multiple memory cells; they are functionally interchangeable.

By embedding these updated LSTM variants within residual block modules, we obtain xLSTM blocks, as detailed in Section~\ref{sec:xLSTMblock}. Architectures built by stacking these xLSTM blocks in a residual manner are referred to as xLSTM architectures (see Section~\ref{sec:xLSTMarchitecure}). Figure~\ref{fig:xlstm_sketch} illustrates the structure of the xLSTM architecture and its constituent parts.

\subsection{Review of the Long Short-Term Memory}
\label{sec:orgLSTM}

The foundational concept behind LSTM~\citep{Hochreiter:91,Hochreiter:97nips,Hochreiter:97} involves utilizing a scalar memory cell as the core unit for both computation and storage. This architecture effectively counters the vanishing gradient problem~\citep{Hochreiter:91,Hochreiter:00book} via the so-called constant error carousel, which pertains to the method by which the cell state is updated. Each memory cell operates under the control of three distinct gating mechanisms: input, output, and forget gates, with the forget gate being subsequently incorporated as detailed by \citet{Gers:00}. The dynamics governing the LSTM memory cell at each time point $t$ are as follows:

\begin{align}
c_t \ &= \  f_t \ c_{t-1} \ + \ i_t \ z_t &  & &\text{cell state} \\
h_t  \ &= \ o_t \ \tilde{h}_t \ , & \tilde{h}_t \ &= \ \psi \left( c_t\right)
  &\text{hidden state} \\
z_t \ &= \ \varphi \left( \tilde{z}_t \right) \ , 
  &\tilde{z}_t \ &=  \ \mathbf{w}^\top_{z} \ \mathbf{x}_t \ + \
  r_{z}  h_{t-1} \ + \  b_{z} \ \
  &\text{cell input} \\
i_t \ &= \ \sigma \left( \tilde{i}_t  \right) \ , 
  &\tilde{i}_t \ &= \ \mathbf{w}^\top_{i} \ \mathbf{x}_t \ + \
  r_{i}  \ h_{t-1} \ + \  b_{i} \ \
  &\text{input gate} \\
f_t \ &= \ \sigma \left(  \tilde{f}_t \right) \ , 
  &\tilde{f}_t \ &= \ \mathbf{w}^\top_{f} \ \mathbf{x}_t  \ + \
  r_{f}  \ h_{t-1} \ + \  b_{f} \ \
  &\text{forget gate} \\
o_t \ &= \ \sigma \left( \tilde{o}_t \right) \ , 
  &\tilde{o}_t  \ &= \ \mathbf{w}^\top_{o} \ \mathbf{x}_t \ + \
  r_{o}  \ h_{t-1} \ + \  b_{o} \ \
  &\text{output gate} 
\end{align}

The vectors $\mathbf{w}_{z}$, $\mathbf{w}_{i}$, $\mathbf{w}_{f}$, and $\mathbf{w}_{o}$ represent the input weights connecting $\mathbf{x}_t$ to the cell input, as well as the input, forget, and output gates, in that order. The parameters $r_{z}$, $r_{i}$, $r_{f}$, and $r_{o}$ denote the recurrent weights linking the previous hidden state $h_{t-1}$ to the cell input, input gate, forget gate, and output gate, respectively. Bias parameters associated with these components are given by $b_{z}$, $b_{i}$, $b_{f}$, and $b_{o}$. The functions $\varphi$ and $\psi$ refer to the activation functions for the cell input and hidden state, which are often chosen as $\tanh$. In particular, $\psi$ serves to bound the range of the cell state, preventing it from diverging. All gate activations use the sigmoid function, expressed as $\sigma \left( x \right)= 1/(1+\exp(-x))$. Subsequent developments have introduced the use of several scalar memory cells $\mathbf{c}_t \in \mathbb{R}^{d}$ grouped into a vector, enabling the implementation of recurrent weight matrices $\mathbf{R} \in \mathbb{R}^{d \times d}$ that blend outputs across memory cells~\citep{Greff:15}; further explanation is available in Appendix~\ref{sec:apporgLSTM}. Ablation analyses indicate that each part of the memory cell architecture plays an essential role~\citep{Greff:15}.

\subsection{sLSTM}
\label{sec:sLSTM}

To enable LSTMs to modify storage choices dynamically, we incorporate exponential gates (depicted in red) alongside mechanisms for normalization and stabilization. Specifically, we allow the input and forget gates to utilize exponential activation functions. As part of the normalization process, we define a normalizer state that accumulates the sum of the input gate multiplied by every subsequent forget gate.

The scalar sLSTM forward pass is:

\begin{align}
\label{eq:slstmforward}
c_t \ &= \  f_t \ c_{t-1} \ + \ i_t \ z_t & &
 &\text{cell state} \\
n_t \ &= \  f_t \ n_{t-1} \ + \ i_t  & &
 &\text{normalizer state} \\
h_t \ &= \ o_t \ \tilde{h}_t \ ,
  &\tilde{h}_t \ &= \ c_t / n_t
&\text{hidden state} \\
z_t \ &= \ \varphi \left( \tilde{z}_t \right) \ , 
  &\tilde{z}_t \ &=  \ \mathbf{w}^\top_{z} \ \mathbf{x}_t \ + \
  r_{z}  h_{t-1} \ + \  b_{z}
  &\text{cell input} \\
i_t \ &= \ \!\exp \left( \tilde{i}_t  \right) \ , 
  &\tilde{i}_t \ &= \ \mathbf{w}^\top_{i} \ \mathbf{x}_t \ + \
  r_{i}  \ h_{t-1} \ + \  b_{i}
  &\text{input gate} \\
f_t \ &= \ \sigma \left(  \tilde{f}_t \right) \ \text{OR} \
\!\exp \left(  \tilde{f}_t \right) \ ,
  &\tilde{f}_t \ &= \ \mathbf{w}^\top_{f} \ \mathbf{x}_t  \ + \
  r_{f}  \ h_{t-1} \ + \  b_{f}
  &\text{forget gate} \\
o_t \ &= \ \sigma \left( \tilde{o}_t \right) \ , 
  &\tilde{o}_t  \ &= \ \mathbf{w}^\top_{o} \ \mathbf{x}_t \ + \
  r_{o}  \ h_{t-1} \ + \  b_{o}
  &\text{output gate}
\end{align}

The original LSTM gating mechanisms—specifically, input- and hidden-dependent gates combined with a bias term—are adapted for these novel architectures. Because exponential activation functions may generate excessively high outputs that risk numerical overflow, we introduce an extra state, \mbox{$m_t$~\citep{Milakov:18arxiv}}, to provide stabilization for the gates:

\begin{align}
\label{eq:slstmstabil}
m_t \ &= \ \max \left( \log ( f_t ) + m_{t-1} , \log ( i_t ) \right) &\text{stabilizer state} \\
i'_t \ &= \ \exp \left( \log \left( i_t \right) - m_t \right) \ = \ \exp \left( \tilde{i}_t - m_t \right) \ \, 
  &\text{stabil. input gate} \\
f'_t \ &= \ \exp \left( \log \left( f_t \right) + m_{t-1} - m_t \right) \ \, 
  &\text{stabil. forget gate}
\end{align}

As demonstrated in Appendix~\ref{sec:appsLSTM}, substituting $f_t$ with $f'_t$ and $i_t$ with $i'_t$ during the forward pass leaves both the network's overall output and the gradients of the loss function with respect to the parameters unaltered.

\paragraph{New Memory Mixing.} Similar to the original LSTM, sLSTM can utilize several memory cells (refer to Appendix~\ref{sec:appsLSTM}). Employing multiple memory cells permits the blending of memory content through recurrent connections $\mathbf{R}_{\mathbf{z}}$, $\mathbf{R}_{\mathbf{i}}$, $\mathbf{R}_{\mathbf{f}}$, and $\mathbf{R}_{\mathbf{o}}$, which transfer information from the hidden state vector $\mathbf{h}$ to both the memory cell input $\mathbf{z}$ and the gates $\mathbf{i}$, $\mathbf{f}$, and $\mathbf{o}$. The process of memory mixing now incorporates a novel factor: the use of exponential gating. In the redesigned sLSTM, each head may support internal memory mixing among its own memory cells, while such interaction does not occur between different heads. The combination of multi-headed structure and exponential gating presents a fresh mechanism for memory mixing within the sLSTM architecture.

\subsection{mLSTM}
\label{sec:mLSTM}

In order to improve the storage capacity of LSTMs, we substitute the conventional scalar memory cell $c \in \mathbb{R}$ with a matrix form $\mathbf{C} \in \mathbb{R}^{d \times d}$. This adjustment means that memory retrieval now utilizes matrix multiplication. At each timestep $t$, the model is designed to memorize a pair of vectors: a key $\mathbf{k}_t \in \mathbb{R}^d$ and a value $\mathbf{v}_t \in \mathbb{R}^d$, following the notation established in the Transformer literature. Subsequently, at time $t+\tau$, the stored value $\mathbf{v}_t$ should be recovered through the use of a query vector $\mathbf{q}_{t+\tau} \in \mathbb{R}^d$. This memory storage and retrieval protocol corresponds to the framework of Bidirectional Associative Memories (BAMs) \citep{Kohonen:72,Anderson:72,Nakano:72,Anderson:77}.

The covariance update rule~\citep{Sejnowski:77,Dayan:91} for storing a key-value pair is

\begin{align}
\mathbf{C}_t \ &= \ \mathbf{C}_{t-1} \ + \ \mathbf{v}_{t} \ \mathbf{k}_t^\top \ .
\end{align}

We apply a layer normalization to the inputs prior to their projection onto keys and values, resulting in vectors with zero mean. This covariance update strategy has been shown to be optimal~\citep{Dayan:91} for maximizing the distinguishability among retrieved binary vectors, which corresponds to maximizing the signal-to-noise ratio. Restricting memory retrieval to pairwise associations and accepting quadratic computational demands, as seen in attention mechanisms~\citep{Krotov:16,Krotov:17,Ramsauer:21}, can yield even greater separability. The covariance update mechanism is identical to that used in Fast Weight Programmers \citep{Schmidhuber:92ncfastweights,Schlag:21}, which were subsequently augmented to incorporate a constant decay factor applied to $\mathbf{C}_{t-1}$ and a fixed learning rate factor for 
$\mathbf{v}_{t} \mathbf{k}_t ^\top$~\citep{Ba:16}. Building upon this foundation, we embed the covariance update rule within the LSTM architecture, mapping the forget gate to the decay rate, the input gate to the learning rate, and employing the output gate to scale the retrieved representation.

In this matrix memory framework, the normalizer state is constructed as a linear combination of key vectors, with each vector scaled by its corresponding input gate and the multiplicative product of subsequent forget gates. This normalizer state thus tracks the cumulative influence of the gates. Because the dot product between the query and the normalizer state may approach zero, the absolute value of this product is utilized and is further constrained by a lower threshold (commonly set to 1.0) in alignment with prior work~\citep{Sun:23arxiv}. The resulting mLSTM forward computation is:

\begin{align}
\label{eq:mlstm_recurrent_begin}
\mathbf{C}_t \ &= \  f_t \ \mathbf{C}_{t-1} \ + \ 
  i_t \ \mathbf{v}_t \ \mathbf{k}_t^\top & &  &\text{cell state} \\
\mathbf{n}_t \ &= \  f_t \ \mathbf{n}_{t-1} \ + \ 
  i_t \ \mathbf{k}_t & &  &\text{normalizer state} \\
\mathbf{h}_t  \ &= \ \mathbf{o}_t \ \odot \ \tilde{\mathbf{h}}_t
\ , \qquad \qquad 
\tilde{\mathbf{h}}_t \ = \   \frac{\mathbf{C}_t \mathbf{q}_t}{\max \left\{ |\mathbf{n}_t^\top \mathbf{q}_t|, 1 \right\} }
   &&&\text{hidden state} \\
\mathbf{q}_t \ &= \ \mathbf{W}_q \ \mathbf{x}_t \ + \ \mathbf{b}_q  & &  &\text{query input} \\
\mathbf{k}_t \ &= \ \frac{1}{\sqrt{d}} \mathbf{W}_k \ \mathbf{x}_t \ + \ \mathbf{b}_k  & &  &\text{key input} \\
\mathbf{v}_t \ &= \ \mathbf{W}_v \ \mathbf{x}_t \ + \ \mathbf{b}_v  & &  &\text{value input} \\
i_t \ &= \  \exp\!\left( \tilde{i}_t \right) \ , 
 \qquad \qquad \quad
  \tilde{i}_t \ = \ \mathbf{w}^\top_{i} \ \mathbf{x}_t \ + \  b_{i} \ \
  &&&\text{input gate} \\
f_t \ &= \ \sigma \!  \left(  \tilde{f}_t \right) \ \text{OR} \ \exp\!\left(  \tilde{f}_t \right) , \ \, 
\tilde{f}_t \ = \ \mathbf{w}^\top_{f} \ \mathbf{x}_t  \ + \
  b_{f}
  &&& \text{forget gate} \\
\label{eq:mlstm_recurrent_end}
\mathbf{o}_t \ &= \ \sigma \left( \tilde{\mathbf{o}}_t \right) \ , \qquad \qquad \quad
  \tilde{\mathbf{o}}_t  \ = \ \mathbf{W}_{\mathbf{o}} \ \mathbf{x}_t \ + \
  \mathbf{b}_{\mathbf{o}}
  &&&\text{output gate} 
\end{align}

Similar to the original LSTM, mLSTM is capable of accommodating several memory cells. Within mLSTM, having multiple heads is functionally identical to having multiple cells because memory elements remain independent without any mixing. To ensure the stability of mLSTM's exponential gating mechanisms, we employ stabilization methods analogous to those used for sLSTM, as described in Equation~\ref{eq:slstmstabil}. Owing to the absence of memory mixing in mLSTM, its recurrence relationship lends itself to a parallelized implementation. Additional explanations can be found in Appendix~\ref{sec:appmLSTM}.

\subsection{xLSTM Architecture}
\label{sec:xLSTMarch}

\paragraph{xLSTM Blocks.}
\label{sec:xLSTMblock}
The role of an xLSTM block is to non-linearly encode historical information within a high-dimensional feature space, thereby facilitating more effective disambiguation of distinct contextual histories. Successfully distinguishing between different histories is essential for the accurate prediction of subsequent sequence elements, such as the upcoming token. Our design is guided by Cover’s Theorem~\citep{Cover:65}, which suggests that patterns embedded non-linearly in higher-dimensional spaces are more amenable to linear separation than those residing in the original input space. We explore two variants of residual block architectures: (i) The first, analogous to those in Transformer models, employs a post up-projection strategy where the historical information is initially transformed non-linearly within the original dimensionality, followed by a linear projection to an expanded space, the application of a non-linear activation, and a linear projection returning the representation to its original space; this configuration corresponds to the left panel in Figure~\ref{fig:Backbones} and the third column of Figure~\ref{fig:xlstm_sketch}. A comprehensive illustration can be found in Figure~\ref{fig:appsLSTM_detailed} of the appendix. (ii) The second variant, inspired by State Space Models, utilizes a pre up-projection scheme where the input is first linearly projected into a higher-dimensional space,

The approach first performs a non-linear compression of past information within a high-dimensional space, followed by a linear transformation back to the initial feature space. When an xLSTM block incorporates an sLSTM, we employ the post up-projection block. Conversely, when an xLSTM block incorporates an mLSTM, the pre up-projection block is utilized due to the expanded memory capacity afforded by the high-dimensional representation. For additional clarification, consult the left panel of Figure~\ref{fig:Backbones}, the third column of Figure~\ref{fig:xlstm_sketch}, or Figure~\ref{fig:appsLSTM_detailed} in the appendix.

\paragraph{xLSTM Architecture. }
\label{sec:xLSTMarchitecure}
The xLSTM model is formed by stacking multiple building blocks with residual connections~\citep{Srivastava:15,He:16}. For its backbone, we adopt the widely used pre-LayerNorm~\citep{Ba:16b} residual architecture, aligning with practices in state-of-the-art Large Language Models. This configuration is illustrated in the final column of Figure~\ref{fig:xlstm_sketch}.

\subsection{Memory and Speed Considerations}

In contrast to Transformers, xLSTM networks feature linear computational complexity and maintain fixed memory requirements regardless of the length of the sequence. Due to the compressive nature of xLSTM memory, these models are highly appropriate for use in industrial contexts and for deployment on edge devices.

The mLSTM's memory operates without the need for additional parameters; however, it incurs significant computational cost due to the $d \times d$ size of both its memory and the update operations. This design necessitates a balance between increased memory capacity and the resulting computational demands. Despite this, the operations can be parallelized efficiently on GPUs, which minimizes their actual runtime impact.

Although mLSTM can be parallelized similarly to approaches such as FlashAttention~\citep{Dao:22,Dao:23} or GLA~\citep{Yang:23arxiv}, sLSTM does not support parallelization owing to the memory mixing involved in its hidden-to-hidden interactions. Nevertheless, we have engineered an efficient CUDA-based solution that employs GPU memory optimizations down to the register level, resulting in sLSTM execution that is generally under twice as slow as mLSTM.

\section{Related Work}
\label{sec:related}

\paragraph{Linear Attention.} Numerous approaches have been proposed to address the quadratic scaling of Transformer attention with respect to context length, aiming to achieve linear complexity. The Synthesizer generates synthetic attention weights, thus removing direct token-to-token dependencies~\citep{Tay:20arxiv}. Linformer implements self-attention using a low-rank projection, providing a linear approximation~\citep{Wang:20arxiv}. The Linear Transformer modifies the attention function itself to achieve linearization~\citep{Katharopoulos:20}. Performer attains a linear approximation to the softmax-based attention by employing a method based on positive orthogonal random features~\citep{Choromanski:21}. Furthermore, in Structured Global Convolution (SGConv)~\citep{Li:22} and the Hyena Hierarchy~\citep{Poli:23}, attention is substituted by fast long convolutional operations.

\paragraph{State Space Models.} In recent times, State Space Models (SSMs) have attracted significant interest due to their linear complexity with respect to sequence length and their competitive results when benchmarked against Transformers. Among the earliest in this category was the Structured State Space sequence model (S4)~\citep{Gu:21}. This was succeeded by a variety of extensions and alternative architectures, including the Diagonal State Space (DSS) model~\citep{Gupta:22}, Gated State Space (GSS) models~\citep{Mehta:22}, S5 model~\citep{Smith:22}, Bidirectional Gated SSM (BiGS)~\citep{Wang:22}, H3 model~\citep{Fu:23}, and Mamba~\citep{Gu:24arxiv}.

\paragraph{Recurrent Neural Networks.} Recurrent Neural Networks (RNNs) have emerged as alternatives to Transformer and attention mechanisms, primarily because their computational complexity scales linearly with the input sequence length. Implementations utilizing Deep Linear Recurrent Units (LRUs) have demonstrated strong performance in language modeling tasks~\citep{Orvieto:23, De:24arxiv}. Other architectures, such as the Hierarchically Gated Linear RNN (HGRN)~\citep{Qin:23} and its successor HGRN2~\citep{Qin:24arxiv}, have reported similar success. Among the most notable RNN-based frameworks for large language models is RWKV~\citep{Peng:23arxivshort,Peng:24arxivshort}, which achieves results comparable to those of Transformer models.

\paragraph{Gating.}
A central concept introduced by LSTM is gating, a mechanism that has subsequently been independently discovered and adapted in numerous recent methods. Gating mechanisms have been incorporated into architectures such as HGRN~\citep{Qin:23}, HGRN2~\citep{Qin:24arxiv}, Gated Linear Attention (GLA)~\citep{Yang:23arxiv}, Gated State Space (GSS) models~\citep{Mehta:22}, Bidirectional Gated SSM (BiGS)~\citep{Wang:22}, Moving Average Equipped Gated Attention (MEGA)~\citep{Ma:22}, RWKV~\citep{Peng:23arxivshort}, as well as Mamba~\citep{Gu:24arxiv}.

\paragraph{Covariance Update Rule.} To improve memory capacity, we integrated a matrix memory governed by a covariance update rule into the mLSTM cell. Various approaches employing similar update strategies include Fast Weight Programmers \citep{Schmidhuber:92ncfastweights,Schlag:21}, RWKV-5 and RWKV-6~\citep{Peng:24arxivshort}, Retention~\citep{Sun:23arxiv}, Linear Transformer~\citep{Katharopoulos:20}, and HGRN2~\citep{Qin:24arxiv}.

\paragraph{Most Related.} The architectures most analogous to xLSTM are Retention~\citep{Sun:23arxiv}, RWKV~\citep{Peng:23arxivshort,Peng:24arxivshort}, and HGRN2~\citep{Qin:24arxiv}, as these models implement either a matrix memory mechanism, gating, or both. Nevertheless, unlike the new sLSTM, they lack the capability for memory mixing. This feature—memory mixing—is essential for effective state tracking, which enables LSTMs to handle tasks that State Space Models (SSMs) and Transformers~\citep{Merrill:24,Deletang:23} are less equipped for, thus rendering LSTMs more expressive. Accurate state tracking is necessary for tasks such as code execution or following entities throughout extended narratives.

\paragraph{Residually Stacking Architectures.} Consistent with nearly all state-of-the-art large-scale deep learning models, xLSTM models employ a design in which their constituent modules are assembled through residual stacking~\citep{Srivastava:15,He:16}.

This approach has been foundational in the development of deep convolutional neural networks~\citep{He:16} as well as Transformers~\citep{Vaswani:17}. The Transformer architecture, in turn, underpins the success of Large Language Models (LLMs) such as GPT-3~\citep{Brown:20short}, ChatGPT~\citep{Schulman:22short}, GPT-4~\citep{Achiam:23gpt4short}, Megatron-LM~\citep{Shoeybi:19}, Gopher~\citep{Rae:21short}, ERNIE 3.0 Titan~\citep{Wang:21arxivshort}, GLaM~\citep{Du:21glamshort}, Chinese M6~\citep{Lin:21}, AlexaTM 20B for multilingual applications~\citep{Soltan:22}, OPT~\citep{Zhang:22opt}, Chinchilla~\citep{Hoffmann:22short}, BLOOM~\citep{Scao:22arxivshort}, GLM-130B~\citep{Zeng:22arxivshort}, LaMDA~\citep{Thoppilan:22short}, PaLM~\citep{Chowdhery:22short}, Llama~\citep{Touvron:23arxiv}, and Gemini~\citep{GeminiTeam:23arxiv,Reid:24arxiv}.

\section{Experiments}
\label{sec:Exp}

This section presents an empirical study of xLSTM, benchmarking its performance against leading approaches within the field of language modeling. In Section~\ref{sec:ExpSyntethic}, we examine xLSTM's distinctive abilities using carefully constructed synthetic benchmarks. Section~\ref{sec:ExpComparison} features a comparison of validation perplexities for state-of-the-art language models, all trained on a dataset of 15B tokens sourced from SlimPajama~\citep{Soboleva:23}. Here, we additionally conduct ablation experiments to better understand the contributions of various components within xLSTM. We further analyze the scaling trends of different models in alignment with findings from \citet{Kaplan:20} and \citet{Brown:20short}. Section~\ref{sec:ExpLanguage} extends the investigation through an extensive language modeling analysis, where both xLSTM and top competitors from Section~\ref{sec:ExpComparison} are trained with a larger corpus of 300B SlimPajama tokens~\citep{Soboleva:23}. This part begins by exploring each method’s capacity to extrapolate over longer textual contexts, evaluates them in terms of validation perplexity and downstream task outcomes~\citep{Sutawika:24}, assesses their effectiveness on 571 text domains drawn from the PALOMA language benchmark~\citep{Magnusson:23arxivshort}, and finally, revisits the question of scaling behavior, this time using an order of magnitude (20$\times$) more training examples.

Throughout our experiments, we denote xLSTM[$a$:$b$] as representing a ratio of $a$ mLSTM-based xLSTM blocks to $b$ sLSTM-based blocks. For instance, the configuration xLSTM[7:1] indicates a structure where, among eight total blocks, seven utilize mLSTM and one utilizes sLSTM. Applied to a typical architecture with 48 total blocks, this notation corresponds to allocating 42 blocks to mLSTM and 6 to sLSTM. Additionally, we consistently incorporate up-projection blocks before the mLSTM modules and after the sLSTM modules across all experimental settings.

\subsection{Synthetic Tasks and Long Range Arena}
\label{sec:ExpSyntethic}
We begin by evaluating the capabilities of xLSTM’s novel exponential gating combined with memory mixing on formal language tasks~\citep{Deletang:23}. Subsequently, we investigate the performance of xLSTM’s newly introduced matrix memory within the context of the Multi-Query Associative Recall task~\citep{Arora:23arxiv}. Lastly, we examine how xLSTM handles long sequence processing by benchmarking it on the Long Range Arena~\citep{Tay:21}.

\paragraph{Evaluation of xLSTM's Exponential Gating with Memory Mixing.} We evaluate the new exponential gating mechanism with memory mixing in xLSTM, which is hypothesized to address challenges in state tracking tasks~\citep{Merrill:24, Merrill:23}. Building on the formal language benchmarks from \citet{Deletang:23}, we modify and expand these tasks to support multi-length training, allowing us to investigate models' ability to extrapolate to sequences of unseen lengths. For comprehensive task descriptions and extended empirical results, refer to Appendix~\ref{sec:appExpSynthetic-formalLang}. We conduct comparisons between xLSTM and a range of baseline architectures, including Transformer models, State Space Models, and various Recurrent Neural Network configurations. We measure model performance exclusively on task-specific token outputs, reporting accuracies normalized between 0 (chance performance) and 1 (ideal performance). Our study encompasses 2-block model variants from several approaches: xLSTM[0:1] (comprising only sLSTM), xLSTM[1:0] (containing only mLSTM), xLSTM[1:1], as well as Llama, Mamba, RWKV, Retention, Hyena, LSTM, and Transformer blocks integrating LSTM (LSTM (Block)). Experimental results are presented in Figure~\ref{fig:formal-main}. Notably, models such as Transformers and State Space Models that lack memory mixing—thereby lacking state tracking ability—are unable to successfully handle tasks involving regular grammars, such as the parity task. This observation is consistent with prior findings indicating that the computational expressivity of Transformers and State Space Models is inherently inferior to that of RNNs~\citep{Merrill:24, Merrill:23, Deletang:23}.

\paragraph{Test of xLSTM's Memory Capacities on Associative Recall Tasks.} In this study, we evaluate the matrix-based memory mechanism of xLSTM on the Multi-Query Associative Recall task~\citep{Arora:23arxiv}, a benchmark designed to probe models' associative memory limits: Each input sequence comprises key-value pairs sampled randomly from an extensive vocabulary, which the model is required to memorize and recall upon query. To provide a more rigorous evaluation than the original setup, we escalate the number of key-value pairs as high as 256, and extend the context window up to 2048 tokens. This expanded protocol allows for a more comprehensive assessment of the comparative memory capacities among various models. Our evaluation covers 2-block configurations of Llama, Mamba, RWKV-5, RWKV-6, xLSTM[1:1], and xLSTM[1:0], measuring their accuracy in retrieving the correct key-value pairs. Due to their exponential scaling of memory with respect to the representational dimension~\citep{Ramsauer:21}, Transformer-based models such as Llama serve as the upper benchmark for this task. Performance outcomes are depicted in Figure~\ref{fig:mqar-main}. Notably, xLSTM[1:1] leads the field among non-Transformer models, maintaining this edge even in smaller-scale settings. Contrary to expectations, the inclusion of the sLSTM block does not negatively affect memory capacity; in fact, its integration appears beneficial, particularly as demonstrated in the most challenging case involving 256 pairs. Further analyses can be found in Appendix~\ref{sec:appExpSynthetic-mqar}, including extrapolation studies that suggest xLSTM's augmented memory design facilitates generalization to context lengths beyond those encountered in training.

\paragraph{Test of xLSTM's Long Context Capabilities on Long Range Arena.} To evaluate the ability of xLSTM to manage extended sequences and broader contexts, we benchmark various approaches using the Long Range Arena~\citep{Tay:21}. Across all tested tasks, xLSTM reliably achieves high levels of performance, indicating that its architecture is notably effective for a range of long context challenges.

For more details, see Appendix~\ref{sec:appExpSynthetic-lra}.

\subsection{Method Comparison and Ablation Study}
\label{sec:ExpComparison}

In pursuit of our central research question—namely, evaluating the performance of our novel LSTM variants when scaled for language modeling tasks—we conduct experiments where xLSTMs, along with Transformers, State Space Models, and additional approaches, are trained on a dataset comprising 15B tokens from SlimPajama, all within an auto-regressive framework. We assess these models based on their validation set performance and additionally carry out ablation analyses specifically targeting xLSTMs.

\paragraph{Comparing xLSTM to Other Methods.} To perform a fair comparison, we train each model on 15B tokens from the SlimPajama dataset~\citep{Soboleva:23} and assess their performance using perplexity scores on the validation split. The following methods are included in the evaluation: our proposed xLSTM, GPT-3 (a Transformer-based model)~\citep{Brown:20short}, Llama (Transformer)~\citep{Touvron:23arxiv}, H3 (a State Space Model or SSM)~\citep{Fu:23}, Mamba (SSM)~\citep{Gu:24arxiv}, RWKV-4 (RNN)~\citep{Peng:23arxivshort}, RWKV-5 and RWKV-6 (RNNs)~\citep{Peng:24arxivshort}, GLA (a linear Transformer)~\citep{Yang:23arxiv}, HGRN and HGRN2 (RNNs)~\citep{Qin:23,Qin:24arxiv}, RetNet (linear Transformer)~\citep{Sun:23arxiv}, Hyena (linear Transformer)~\citep{Poli:23}, as well as xLSTM variants xLSTM[1:0] and xLSTM[7:1] (definitions in Section~\ref{sec:xlstm_blocks_notation}). 
All models utilize mixed precision for training; however, RWKV-5, RWKV-6, GLA, and HGRN2 do not rely on PyTorch's automated mixed-precision functionality (see Appendix Section~\ref{sec:appGenTrainProc} for additional implementation notes). 
The approaches are grouped into: (a) Transformers, (b) State Space Models (SSMs), and (c) Recurrent Neural Networks (RNNs) and linear Transformers. Linear Transformers replace the traditional Transformer attention with linear alternatives. To standardize comparisons, the models are designed to align with a 350M-parameter GPT-3 configuration, featuring an embedding dimension of 1024 and 24 residual blocks. Only GPT-3 shares its token and output embedding matrices, resulting in a reduced parameter count. According to Table~\ref{tab:spaj15b_model_results}, xLSTM achieves the lowest validation perplexity, surpassing all prior baselines. Further experimental details are available in Appendix~\ref{sec:appExpComparison}. The scaling results displayed in Figure~\ref{fig:temp_sclaw15B} suggest that xLSTM is likely to maintain its advantages even as model size increases.

\paragraph{Ablation Studies.}
As shown in Table~\ref{tab:spaj15b_model_results} and Figure~\ref{fig:temp_sclaw15B}, the xLSTM model delivers strong language modeling performance when trained on the 15B-token SlimPajama dataset. To systematically assess the contributions of each component in xLSTM, we progressively transition a standard LSTM model into the xLSTM variant. The ablation is performed by first inserting LSTM layers within pre-LayerNorm residual architectures. Next, the model is further modified with the inclusion of a post up-projection block. In the final stage, we incorporate both exponential gating and the matrix memory mechanism.

The results are shown in Table~\ref{tab:ablstudies} (top). 

Our ablation experiments reveal that both exponential gating and matrix memory substantially contribute to the observed performance gains. Recognizing the central role of gating within RNNs and State Space Models, we further investigate the effects of varying gating approaches. As shown in Table~\ref{tab:ablstudies} (bottom), enabling each gate to be both learnable and dependent on the input leads to stepwise improvements. Further analysis concerning the distinct backbone modules is provided in Appendix~\ref{sec:appAblDetails}.

\subsection{xLSTM as Large Language Model}
\label{sec:ExpLanguage}

This research concludes with extensive language modeling experiments aimed at assessing the efficacy of xLSTM as a large language model. To this end, we scale up the training data, utilizing 300B tokens sourced from SlimPajama. Notably, this quantity of data aligns with that employed in prior works such as Mamba~\citep{Gu:24arxiv} and Griffin~\citep{De:24arxiv}.

We perform a comparative analysis of xLSTM against RWKV-4, Llama, and Mamba, each exemplifying one of the methodological categories detailed in Section~\ref{sec:ExpComparison}. RWKV-4 is chosen as the RNN exemplar, as optimal training precision for RWKV-5, RWKV-6, and HGRN2 was only determined following the initiation of experiments using 300B tokens (refer to Appendix~\ref{sec:appGenTrainProc}). A range of model scales is trained (125M, 350M, 760M, and 1.3B parameters), with evaluations including tests for length generalization and assessments on validation datasets. Additionally, the models' downstream task performance is analyzed, their language modeling abilities are benchmarked across 571 text domains from PALOMA, and their scaling law characteristics are explored.

\paragraph{Sequence Length Extrapolation.}
\label{sec:spaj300B_ctx_extrapolation}
To begin, we evaluate sequence length extrapolation capabilities for xLSTM, RWKV-4, Llama, and Mamba models at the 1.3B parameter scale. Training for all models is performed with a maximum context length of 2048, after which they are assessed at extended context lengths reaching up to 16384. Results are presented in Figure~\ref{fig:spaj300B_seq_extrapolation}. Notably, xLSTM models exhibit consistently low perplexity across longer context windows, outperforming the alternative approaches in this regard.

\paragraph{Validation Perplexity and Downstream Tasks.}
\label{sec:spaj_downstream_eval}
Next, we assess all model variants—xLSTM, RWKV-4, Llama, and Mamba—across various model sizes by measuring their next token prediction capabilities on the SlimPajama validation dataset, as well as their effectiveness on downstream benchmarks that test for common sense reasoning skills.

The third column in Table~\ref{tab:lmeval} presents the perplexity scores on the validation set for each method evaluated. Notably, xLSTM[1:0] and xLSTM[7:1] achieve the lowest validation set perplexities across all model sizes examined, making them the top-performing models on this metric. Additional columns in Table~\ref{tab:lmeval} show results on various downstream benchmarks. Across the majority of tasks and model configurations, xLSTM consistently yields the strongest results; Mamba surpasses xLSTM only on the ARC task in select instances. Further experimental details are provided in Appendix~\ref{sec:appExpLanguage}.

\paragraph{Performance on PALOMA Language Tasks.} Third, we evaluate the next token prediction abilities of xLSTM, RWKV-4, Llama, and Mamba models, across all parameter scales, on the suite of PALOMA language tasks~\citep{Magnusson:23arxivshort}. The primary metric is perplexity, calculated for next token prediction across 571 diverse textual sources, encompassing everything from nytimes.com articles to Reddit communities like r/depression.  
Table~\ref{tab:ppleval} presents the resulting perplexity scores, organized into two categories: language modeling domains (the first seven columns), and a selection of detailed domain benchmarks (the last five columns). Among these tasks, the xLSTM[1:0] configuration achieves superior performance compared to the xLSTM[7:1] variant.

xLSTM[1:0] achieves lower perplexity compared to Mamba in 568 of 571 text domains (99.5\%), outperforms Llama in 486 of 571 domains (85.1\%), and surpasses RWKV-4 in 570 of 571 cases (99.8\%). Additional results can be found in Appendix~\ref{sec:appExpLanguage}.

\paragraph{Scaling Laws.} Next, we investigate the adherence to power-law scaling trends, which serve as a basis for predicting how performance will change with increasing model size~\citep{Kaplan:20,Brown:20short}. The scaling trends for the different architectures are illustrated in Figure~\ref{fig:temp_sclaw300B}. While all models exhibit comparable scaling patterns, they differ by additive constants. Among the evaluated architectures, RWKV-4 yields the lowest performance, with Llama and Mamba achieving progressively better results, and xLSTM surpassing Mamba by a margin similar to the gap between Mamba and Llama. These observations suggest that, as model size increases, xLSTM is likely to sustain an advantage over both Transformer-based and State-Space-based counterparts.

\textbf{Generation Times and Maximal Throughput.} We evaluate the time required for text generation in Figure~\ref{fig:inference_time_generation} and assess the maximum achievable throughput in Figure~\ref{fig:inference_time_decoding_throughput} (left) for xLSTM models at the 1.3B parameter scale. Our analysis includes comparisons with equivalently sized Mamba, Llama, and RWKV models from HuggingFace, and for Llama we utilize a static key-value cache configuration. During these experiments, neither full cache compilation for the Transformer model nor \texttt{torch.compile} support for the Mamba architecture were functional. Text generation is benchmarked with each model using a batch size of 1 and a pre-fill value of 16, a setting that particularly benefits the Transformer architecture.

According to Figure~\ref{fig:inference_time_generation}, the computational cost of text generation for xLSTM and the recurrent models Mamba and RWKV-4 scales linearly with sequence length, in contrast to Llama, which demonstrates quadratic scaling behavior. For the assessment of decoding throughput, various batch sizes and pre-fill configurations are considered, specifically for the Llama baseline.

As depicted in Figure~\ref{fig:inference_time_decoding_throughput} (right), xLSTM maintains constant memory usage regardless of batch size, permitting much larger batch processing compared to Llama and resulting in superior throughput.

\section{Limitations}

(i) Unlike mLSTM, the design of sLSTM incorporates memory mixing that inhibits the use of parallelizable operations, precluding efficient parallel execution. Despite this, we engineered a custom CUDA kernel for sLSTM, achieving execution speeds that are currently under twice as slow as our parallel mLSTM implementation. (ii) The CUDA kernels employed for mLSTM have not undergone significant optimization, and as a result, our implementation is approximately four times slower than FlashAttention or the scan operation employed by Mamba. Improved CUDA performance akin to that achieved by FlashAttention could potentially be realized with further optimization efforts. (iii) mLSTM’s use of a matrix memory incurs considerable computational overhead, as it requires operations on $d \times d$ matrices. Nevertheless, the update and readout procedures for the memory are parameter-free and amenable to parallel execution with conventional matrix routines, such that the overall increase in wall time due to the complex memory structure remains modest. (iv) Careful selection is essential when initializing the forget gates. (v) Since the matrix memory’s capacity does not scale with sequence length, utilizing longer sequences could risk saturating the memory at large context lengths. However, no adverse effects have been observed for contexts up to 16k, as described in Section~\ref{sec:spaj300B_ctx_extrapolation}. (vi) The computational demands of large-scale language modeling experiments limited our ability to exhaustively tune either the model design or the hyperparameters, particularly for expansive xLSTM variants. We believe xLSTM can achieve further gains given a thorough optimization process.

\section{Conclusion}
Our investigation has yielded a partial response to the central question: To what extent can language modeling be advanced by scaling LSTM architectures to billions of parameters? The evidence thus far demonstrates that such scaling enables LSTMs to match, at minimum, the capabilities of prevailing approaches such as Transformers and State Space Models. In this work, we introduced xLSTM, an improved variant of LSTM, distinguished by its use of exponential gating with memory mixing and an innovative memory architecture. When benchmarked against leading methods like Transformers and State Space Models, xLSTM consistently exhibits superior or comparable performance in language modeling tasks. The scaling laws derived from our study suggest that with further increases in scale, xLSTM models could emerge as strong alternatives to the existing Large Language Models constructed with Transformer architectures. Beyond language modeling, xLSTM possesses significant promise for application in diverse domains, including Reinforcement Learning, Time Series Prediction, and the simulation of physical processes.

\end{document}